\documentclass[10pt,a4paper]{amsart}
\usepackage{amsmath,amssymb,amsthm,graphicx}
% Retained preamble for the Stacks Project chapter snapshot in this directory.
% The chapter's own preamble.tex is a copy of the shared Stacks Project preamble; the block
% below is sliced MECHANICALLY, line for line, from it. Nothing was added, removed or
% redefined; the wrapper class, the xr-hyper/hyperref wiring and the externaldocument list
% are omitted because the assembler supplies the document class.
%
% lines taken from preamble.tex: 3, 5, 9, 11, 12, 156, 157, 158, ... (51 lines total)
%
% package coverage restored from the chapter's own preamble (2026-10-13): xy, xy-2cell
% package coverage restored from the chapter's own preamble (2026-10-14): verbatim
%   the chapter's preamble.tex loads \usepackage{verbatim} on line 8 and then defines
%   reference/slogan/history as \comment environments; without the package \comment is
%   undefined and any card whose counted spans use one of the three environments fails
%   with "Undefined control sequence" (L44184, lemma-representable-over-space).

\usepackage{verbatim}
\usepackage[all]{xy}
\xyoption{2cell}
\UseAllTwocells

\newenvironment{reference}{\comment}{\endcomment}
\newenvironment{slogan}{\comment}{\endcomment}
\newenvironment{history}{\comment}{\endcomment}
\theoremstyle{plain}
\newtheorem{theorem}[subsection]{Theorem}
\newtheorem{proposition}[subsection]{Proposition}
\newtheorem{lemma}[subsection]{Lemma}
\theoremstyle{definition}
\newtheorem{definition}[subsection]{Definition}
\newtheorem{example}[subsection]{Example}
\newtheorem{exercise}[subsection]{Exercise}
\newtheorem{situation}[subsection]{Situation}
\theoremstyle{remark}
\newtheorem{remark}[subsection]{Remark}
\newtheorem{remarks}[subsection]{Remarks}
\numberwithin{equation}{subsection}
\def\lim{\mathop{\mathrm{lim}}\nolimits}
\def\colim{\mathop{\mathrm{colim}}\nolimits}
\def\Spec{\mathop{\mathrm{Spec}}}
\def\Hom{\mathop{\mathrm{Hom}}\nolimits}
\def\Ext{\mathop{\mathrm{Ext}}\nolimits}
\def\SheafHom{\mathop{\mathcal{H}\!\mathit{om}}\nolimits}
\def\SheafExt{\mathop{\mathcal{E}\!\mathit{xt}}\nolimits}
\def\Sch{\mathit{Sch}}
\def\Mor{\mathop{\mathrm{Mor}}\nolimits}
\def\Ob{\mathop{\mathrm{Ob}}\nolimits}
\def\Sh{\mathop{\mathit{Sh}}\nolimits}
\def\NL{\mathop{N\!L}\nolimits}
\def\CH{\mathop{\mathrm{CH}}\nolimits}
\def\proetale{{pro\text{-}\acute{e}tale}}
\def\etale{{\acute{e}tale}}
\def\QCoh{\mathit{QCoh}}
\def\Ker{\mathop{\mathrm{Ker}}}
\def\Im{\mathop{\mathrm{Im}}}
\def\Coker{\mathop{\mathrm{Coker}}}
\def\Coim{\mathop{\mathrm{Coim}}}
\DeclareMathSymbol{\boxtimes}{\mathbin}{AMSa}{"02}
\def\QCohstack{\mathcal{QC}\!\mathit{oh}}
\def\Cohstack{\mathcal{C}\!\mathit{oh}}
\def\Spacesstack{\mathcal{S}\!\mathit{paces}}
\def\Quotfunctor{\mathrm{Quot}}
\def\Hilbfunctor{\mathrm{Hilb}}
\def\Curvesstack{\mathcal{C}\!\mathit{urves}}
\def\Polarizedstack{\mathcal{P}\!\mathit{olarized}}
\def\Complexesstack{\mathcal{C}\!\mathit{omplexes}}
\def\Pic{\mathop{\mathrm{Pic}}\nolimits}
\def\Picardstack{\mathcal{P}\!\mathit{ic}}
\def\Picardfunctor{\mathrm{Pic}}
\def\Deformationcategory{\mathcal{D}\!\mathit{ef}}
\usepackage{cleveref}
\crefname{originaldefinition}{original definition}{original definitions}
\Crefname{originaldefinition}{Original definition}{Original definitions}
\makeatletter
\newlabel{definition-relative-representable-property}{{025V}{source}{}{sourcepointer-1}{}}
\newlabel{definition-relative-representable-property@cref}{{[originaldefinition][0][]025V}{[1][0][]source}{}{}{}}
\makeatother
\crefname{originallemma}{original lemma}{original lemmas}
\Crefname{originallemma}{Original lemma}{Original lemmas}
\makeatletter
\newlabel{lemma-base-change-representable-transformations-property}{{02WL}{source}{}{sourcepointer-2}{}}
\newlabel{lemma-base-change-representable-transformations-property@cref}{{[originallemma][0][]02WL}{[1][0][]source}{}{}{}}
\makeatother
\crefname{originallemma}{original lemma}{original lemmas}
\Crefname{originallemma}{Original lemma}{Original lemmas}
\makeatletter
\newlabel{lemma-descent-representable-transformations-property}{{03KD}{source}{}{sourcepointer-3}{}}
\newlabel{lemma-descent-representable-transformations-property@cref}{{[originallemma][0][]03KD}{[1][0][]source}{}{}{}}
\makeatother
\title{A property of a representable morphism of algebraic spaces is tested on one etale chart}
\author{The Stacks Project Authors}
\date{stacks-project:spaces.tex:local-snapshot-20261012; GNU FDL 1.2 or later}
\begin{document}
\maketitle
\setcounter{subsection}{1}
\noindent\textit{Editorial note.} All delivered mathematical text is the author's own, copied from the pinned Stacks Project chapter file spaces.tex; no statement or proof step was rewritten, reordered or removed.\par\medskip
\noindent\textit{Editorial note.} This candidate was dropped by the mechanical prefilter because of the phrase 'It is clear from the definitions that ...' at proof line 1982, which opens the easy direction of an if-and-only-if whose converse is proved in full.\par\medskip
\noindent\textit{Editorial note.} The proof contains one commutative diagram; it is part of the counted proof span and is preserved verbatim.\par\medskip
\noindent\textit{Editorial note.} The references of the counted spans are recorded as non-load-bearing navigation pointers; the argument that consumes them stays inside the counted proof span.\par\medskip
\noindent\textit{Editorial note.} Preamble: sources/stacks-spaces/retained-preamble.tex, sliced mechanically from the chapter's own preamble.tex.\par\medskip
\noindent\textit{Allowed standard prerequisites (editorial).}\begin{itemize}
\item Definition of the relative representable property type. Definition definition-relative-representable-property (713-728), used at proof lines 1972-1975. Fixes the properties P to which the statement applies.
\item The two-square diagram comparing V -> U with T' -> T. used at proof lines 1983-1993. Identifies T' -> T as a base change of V -> U.
\item Properties of the type P are local on the base in the fppf topology. used at proof lines 1994-2002. Descends property P along the covering T x\_F U -> T.
\end{itemize}\medskip
\ifdefined\hypertarget\hypertarget{sourcepointer-1}{}\fi
\noindent\textit{Source reference.} Original definition 025V: focused-seeds/AGP/sources/stacks-snapshot-20261011/tags.txt, source lines 2828--2828. Navigation pointer into the Stacks Project; the standard fact it supplies is applied inside the counted proof span and the argument resting on it stays there.\par\medskip
\ifdefined\hypertarget\hypertarget{sourcepointer-2}{}\fi
\noindent\textit{Source reference.} Original lemma 02WL: focused-seeds/AGP/sources/stacks-snapshot-20261011/tags.txt, source lines 3808--3808. Navigation pointer into the Stacks Project; the standard fact it supplies is applied inside the counted proof span and the argument resting on it stays there.\par\medskip
\ifdefined\hypertarget\hypertarget{sourcepointer-3}{}\fi
\noindent\textit{Source reference.} Original lemma 03KD: focused-seeds/AGP/sources/stacks-snapshot-20261011/tags.txt, source lines 4673--4673. Navigation pointer into the Stacks Project; the standard fact it supplies is applied inside the counted proof span and the argument resting on it stays there.\par\medskip
\begin{lemma}
\label{lemma-representable-morphisms-spaces-property}
Let $S$ be a scheme contained in $\Sch_{fppf}$.
Let $F$, $G$ be algebraic spaces over $S$.
Let $G \to F$ be a representable morphism.
Let $U \in \Ob((\Sch/S)_{fppf})$, and $q : U \to F$
surjective and \'etale. Set $V = G \times_F U$.
Finally, let $\mathcal{P}$ be a property of morphisms
of schemes as in Definition \ref{definition-relative-representable-property}.
Then $G \to F$ has property $\mathcal{P}$ if and only if
$V \to U$ has property $\mathcal{P}$.
\end{lemma}

\begin{proof}
(This lemma follows from
Lemmas \ref{lemma-base-change-representable-transformations-property} and
\ref{lemma-descent-representable-transformations-property},
but we give a direct proof here also.)
It is clear from the definitions that if $G \to F$ has property
$\mathcal{P}$, then $V \to U$ has property $\mathcal{P}$.
Conversely, assume $V \to U$ has property $\mathcal{P}$.
Let $T \to F$ be a morphism from a scheme to $F$.
Let $T' = T \times_F G$ which is a scheme since $G \to F$ is
representable. We have to show that $T' \to T$ has property $\mathcal{P}$.
Consider the commutative diagram of schemes
$$
\xymatrix{
V \ar[d] & T \times_F V \ar[d] \ar[l] \ar[r] &
T \times_F G \ar[d] \ar@{=}[r] & T' \\
U & T \times_F U \ar[l] \ar[r] & T
}
$$
where both squares are fibre product squares. Hence we conclude
the middle arrow has property $\mathcal{P}$ as a base change
of $V \to U$. Finally, $\{T \times_F U \to T\}$ is a fppf covering
as it is surjective \'etale, and hence we conclude that
$T' \to T$ has property $\mathcal{P}$ as it is local on the
base in the fppf topology.
\end{proof}

\end{document}
